\scanpage{62}
\mspage{68}

\diagram{\(\begin{array}{lll}
\text{Der 4 (g)}&:&\text{existence des } f_*\\
\text{Der 4 (d)}&:&\text{existence des } f_!\\
\text{\uline{Der} 5 (g)}&&\text{Calcul des fibres de } \deleted{\unsure{R|F}} f_*(F)\\
\text{Der 5 (d)}&&\hspace{1cm}\rule{1.6cm}{0.4pt}\quad \deleted{\unsure{R|F|L}} f_!(F)
\end{array}\)}

L'``ancien axiome Der 5 (g)'' résulte de Der 2, Der 3, Der 4 (g,d), Der 5 (d), grâce au théorème précédent.

Un aspect intéressant du théorème précédent, si on l'envisage comme \unsure{énonçant} l'équivalence de a) et b), c'est que cet énoncé ne fait pas appel à des notions homologiques, mais que sa démonstration y fait appel, sous une forme relativement forte (calcul des fibres de $f_!(F)$). Pour les th. de chang$^{\mathrm t}$ de base pour la $\unsure{\mathbb D|\underline D}$-cohomologie relative i.e. pour les $f_*^{\mathbb D}$), on a besoin aussi constamment du calcul des fibres des $f_*(F)$. Donc il semble qu'on ne peut se passer d'aucun des axiomes précédents, y compris Der 4 (d) et Der 5 (d), même si on ne \inserted{veut} s'intéresser qu'au formalisme des $f^*,f_*$ (aspect cohomologique). Par la suite, on va donc supposer ces axiomes satisfaits.

\uline{Théorème 2} Soit $f:X\unsure{\relbar\joinrel\relbar|\longrightarrow}Y$ dans $\mathrm{Diag}$. Conditions équivalentes

a) $\forall y_0\in\mathrm{Ob}\,Y$, le foncteur canonique \deleted{\unsure{trivial}}

\[ \underset{\text{fibre en }y_0}{X_{y_0}}\xrightarrow{\unsure{\varphi|q}_{\unsure{y|y_0|0}}}X_{/y_0}\qquad (=X\times_Y Y_{/y_0}) \]

est une $\mathbb D$-équivalence costricte.

b) $\forall x_0\in\mathrm{Ob}\,X$, le foncteur induit sur les catégories colocalisées en $x_0$ et en $y_0=f(x_0)$

\[ {}_{x_0\backslash}X\longrightarrow{}_{y_0\backslash}Y \]

est à fibres $\mathbb D$-\unsure{contractiles|contractibles}.
% archive sheet number: 61
